\documentclass[10pt,a4paper]{amsart}
\usepackage[margin=19mm]{geometry}
\usepackage{amsmath,amssymb,mathtools}
\usepackage[scr=rsfs,cal=euler]{mathalfa}
\usepackage{xfrac}
\usepackage[unicode,hidelinks,bookmarks=false]{hyperref}
\newcommand{\ie}{i.e.\ }
\newcommand{\cf}{cf.\ }
\newcommand{\resp}{respectively\ }
\newcommand{\Subsection}{\subsection}
\providecommand{\nobreakdash}{\nobreak\hbox{-}}
\setlength{\emergencystretch}{2em}
\multlinegap0pt
\usepackage{amssymb}
\usepackage{enumerate}
\usepackage{tikz}
\usepackage{xcolor}
\usepackage{dsfont}
\usepackage{float}
\usepackage{bm}
\newtheorem{lem}{Lemma}
\title{Critical coupling thresholds for tilted Kuramoto-Vicsek models with a confining potential}
\author{Benedetta Bertoli, Benjamin D. Goddard, Grigorios A. Pavliotis}
\date{Source: arXiv:2603.18185v1 (CC BY 4.0)}
\begin{document}
\maketitle

\noindent\emph{Source and licence.} Critical coupling thresholds for tilted Kuramoto-Vicsek models with a confining potential. Version arXiv:2603.18185v1 (CC BY 4.0), distributed by arXiv under the Creative Commons Attribution 4.0 International licence, http://creativecommons.org/licenses/by/4.0/, as recorded in the arXiv licence metadata for this exact version and shown on the abstract page https://arxiv.org/abs/2603.18185v1. Full licence notice: \texttt{licenses/arxiv-2603.18185-CC-BY-4.0.txt}.\par\smallskip

\noindent\textit{Editorial scope.} Counted unit: the article's lem lem:uniqueness-homog (source0.tex lines 236--242). The article's complete original proof of it is delivered with it (source0.tex lines 244--258, one \texttt{proof} environment of 15 lines). No mathematics is written, repaired or re-derived: the statement and the proof are the article's own lines, in the article's own notation, and no retained line is substituted or re-flowed.  No further source-local context is needed: every cross-reference of every retained line resolves inside the two delivered spans.  Nothing was omitted from any retained span, so the package carries no omission record.  No retained line contains a citation, so the wrapper carries no bibliography and no bibliography entry.  No retained line contains a figure environment, an image-inclusion command or drawing code, so no image is delivered and the package declares no asset dependency.  Editorial formatting only, in addition: an AMS \texttt{amsart} preamble that restates the article's own declarations and macros used by the retained lines, namely the environments \texttt{lem} on separate counters, together with the standard packages those lines need.  The retained environments print in this document as Lemma 1 in order of appearance, whereas the article prints the same items with the numbering of its own full document.  The complete native source of the article as distributed in the arXiv e-print for this exact version is bundled unchanged as \texttt{sources/196674/source0.tex}.
\begin{lem}\label{lem:uniqueness-homog}
For $\Gamma$ sufficiently large, the map $T:\mathcal{P}\to\mathcal{P}$ defined by
\[
    (T\rho)(\theta) := \frac{1}{Z[\rho]}\,e^{-\Gamma^{-1} U[\rho](\theta)} \int_\theta^{\theta + 2\pi} e^{\Gamma^{-1} U[\rho](\varphi)}\,\mathrm{d}\varphi,
\]
is a contraction on $(\mathcal{P}, \|\cdot\|_{L^1})$, and therefore admits a unique fixed point.
\end{lem}

\begin{proof}
We work throughout on the space $\mathcal{P}$ of $2\pi$-periodic probability densities on $[0,2\pi]$, regarded as a closed subset of $L^1([0,2\pi])$. We first note that $T$ maps $\mathcal{P}$ into itself, which follows from the definition and positivity of $Z[\rho]$. We now establish the contraction. For $\rho,\tilde\rho\in\mathcal{P}$, the two potentials differ only through the Fourier moments $r_{k,c}$ and $r_{k,s}$ , so
\begin{equation}\label{eq:U-Lip}
    \|U[\rho]-U[\tilde\rho]\|_{L^\infty} \leq 2\gamma\sum_{k=1}^n |a_k|\,\|\rho-\tilde\rho\|_{L^1} =: \tilde{A}\,\|\rho-\tilde\rho\|_{L^1},
\end{equation}
using $|\cos(k\theta)|,|\sin(k\theta)|\leq 1$ and $|r_{k,c}[\rho]-r_{k,c}[\tilde\rho]|\leq\|\rho-\tilde\rho\|_{L^1}$. Setting $K[\rho](\theta) := e^{-\Gamma^{-1}U[\rho](\theta)}\int_\theta^{\theta+2\pi}e^{\Gamma^{-1}U[\rho](\varphi)}\,\mathrm{d}\varphi$, the exponential inequality $|e^a - e^b| \leq |a-b|e^{\max(a,b)}$ and the uniform bound $\|U[\rho]\|_{L^\infty}\leq C_0$ for all $\rho\in\mathcal{P}$ together give
\[
    \|K[\rho]-K[\tilde\rho]\|_{L^\infty} \leq C_1\Gamma^{-1}\|U[\rho]-U[\tilde\rho]\|_{L^\infty},
\]
for a constant $C_1 > 0$ depending only on $C_0$. For the normalisation, since $Z[\rho] = \int_0^{2\pi}K[\rho]\,\mathrm{d}\theta$ and $Z[\rho]\geq z_0>0$ uniformly, we have $|Z[\rho]^{-1} - Z[\tilde\rho]^{-1}|\leq z_0^{-2}|Z[\rho]-Z[\tilde\rho]|$. Decomposing $T\rho - T\tilde\rho = Z[\rho]^{-1}(K[\rho]-K[\tilde\rho]) + K[\tilde\rho](Z[\rho]^{-1}-Z[\tilde\rho]^{-1})$ and taking $L^1$-norms, the above estimates give
\[
    \|T\rho - T\tilde\rho\|_{L^1} \leq C_2\,\Gamma^{-1}\|\rho-\tilde\rho\|_{L^1},
\]
for a constant $C_2>0$ independent of $\Gamma$ and of $\rho,\tilde\rho\in\mathcal{P}$. For $\Gamma > C_2$ the map $T$ is therefore a contraction on the complete metric space $(\mathcal{P},\|\cdot\|_{L^1})$, and the Banach fixed-point theorem allows us to conclude uniqueness of the fixed point.
\end{proof}

\end{document}
